% TEMPLATE-MILC-v3.tex - plantilla maestra para todas las guias MILC v3
% Ver scripts/generadores/build-guia-milc.py para la lista completa de
% claves esperadas. El script valida que no queden placeholders sin
% reemplazar antes de escribir el archivo final.

\documentclass[11pt,letterpaper]{article}
\usepackage[margin=1.75cm]{geometry}
\usepackage[table]{xcolor}
\usepackage{fontspec}
\usepackage[spanish]{babel}
\usepackage{tikz}
\usepackage[most]{tcolorbox}
\usepackage{tabularx}
\usepackage{array}
\usepackage{fancyhdr}
\usepackage{titlesec}
\usepackage{ragged2e}
\usepackage{enumitem}
\usepackage{microtype}

\setmainfont{Helvetica Neue}
\setsansfont{Helvetica Neue}
\setlength{\parindent}{0pt}
\setlength{\parskip}{6pt}
\hyphenpenalty=200
\exhyphenpenalty=200
\pretolerance=400
\tolerance=2000
\setlength{\emergencystretch}{1em}
\renewcommand{\arraystretch}{1.25}
\setlist{leftmargin=5mm,itemsep=1mm,topsep=1mm}
\newcolumntype{Y}{>{\RaggedRight\arraybackslash}X}

\definecolor{milcMagenta}{HTML}{E6007E}
\definecolor{milcVino}{HTML}{5A0038}
\definecolor{milcTurquesa}{HTML}{0093A5}
\definecolor{milcVerde}{HTML}{58A923}
\definecolor{milcMostaza}{HTML}{E5B400}
\definecolor{milcOcre}{HTML}{B86600}
\definecolor{milcNegro}{HTML}{241816}
\definecolor{milcGris}{HTML}{F7F7F4}
\definecolor{milcLinea}{HTML}{E8E2DD}
\definecolor{milcGradePrimary}{HTML}{0066FF}
\definecolor{milcGradeDark}{HTML}{003D99}
\definecolor{milcGradeSoft}{HTML}{D6E8FF}
\definecolor{milcGradeAccent}{HTML}{CFE7FF}
\definecolor{milcGradeText}{HTML}{FFFFFF}
\color{milcNegro}

\pagestyle{fancy}
\fancyhf{}
\lhead{\small Tecnología e Informática · Grado 6}
\rhead{\small Guía 1 · MILC v3}
\cfoot{\small\thepage}
\renewcommand{\headrulewidth}{0pt}

\newtcolorbox{titlebox}[1]{
  enhanced,colback=#1,colframe=#1,arc=7mm,boxrule=0pt,
  left=7mm,right=7mm,top=3mm,bottom=3mm,
  before skip=8pt,after skip=8pt,
  fontupper=\sffamily\bfseries\Large\color{white}
}

\newtcolorbox{softbox}[3]{
  enhanced,colback=#2,colframe=#1,borderline west={4pt}{0pt}{#1},
  arc=2mm,boxrule=.4pt,left=4mm,right=4mm,top=3mm,bottom=3mm,
  before skip=6pt,after skip=8pt,title={#3},coltitle=white,
  colbacktitle=#1,fonttitle=\sffamily\bfseries,fontupper=\RaggedRight,
  attach boxed title to top left={xshift=3mm,yshift=-2mm},
  boxed title style={arc=2mm, boxrule=0pt}
}

\newtcolorbox{drawbox}[2]{
  enhanced,colback=white,colframe=#1,arc=4mm,boxrule=1pt,
  left=5mm,right=5mm,top=4mm,bottom=4mm,before skip=8pt,after skip=8pt,
  title={#2},coltitle=white,colbacktitle=#1,fonttitle=\sffamily\bfseries,
  fontupper=\RaggedRight,
  attach boxed title to top left={xshift=4mm,yshift=-2mm},
  boxed title style={arc=3mm, boxrule=0pt}
}

\newtcolorbox{infoband}[1]{
  enhanced,colback=#1!8,colframe=#1!50,arc=2mm,boxrule=.5pt,
  left=4mm,right=4mm,top=2mm,bottom=2mm,before skip=4pt,after skip=4pt,
  fontupper=\small\RaggedRight
}

% Puente narrativo entre secciones (contrato editorial Conéctate).
% Da continuidad: "Acabas de X. Ahora vas a Y porque Z."
% Visualmente: banda izquierda en color del grado, texto en itálica pequeña.
\newtcolorbox{puentebox}{
  enhanced,colback=milcGradeSoft,colframe=milcGradeDark,
  borderline west={3pt}{0pt}{milcGradeDark},
  arc=0mm,boxrule=0pt,left=5mm,right=4mm,top=2.5mm,bottom=2.5mm,
  before skip=4pt,after skip=8pt,
  fontupper=\itshape\small\color{milcNegro}\RaggedRight
}
\newcommand{\puente}[1]{%
  \begin{puentebox}%
    \textcolor{milcGradeDark}{\textbf{\textrightarrow}}\hspace{2mm}#1%
  \end{puentebox}%
}

\newcommand{\linead}[1]{\noindent\color{milcLinea}\rule{#1}{0.5pt}\color{milcNegro}}
\newcommand{\respuesta}[1]{\par\vspace{2mm}\foreach \n in {1,...,#1}{\noindent\color{milcLinea}\rule{\linewidth}{0.5pt}\par\vspace{5mm}}\color{milcNegro}}
\newcommand{\checkbox}{\textcolor{milcGradeDark}{\fbox{\phantom{x}}}\hspace{5pt}}

\newcommand{\phasecard}[4]{
\begin{tcolorbox}[enhanced,colback=#3,colframe=#2,arc=3mm,boxrule=.5pt,height=3.05cm,valign=center,halign=center,left=1.5mm,right=1.5mm,top=2mm,bottom=2mm]
{\sffamily\bfseries\fontsize{22}{24}\selectfont\textcolor{#2}{#1}}\par
{\sffamily\bfseries\small #4}
\end{tcolorbox}
}

\begin{document}
\thispagestyle{empty}
\begin{tikzpicture}[remember picture,overlay]
  \fill[white] (current page.south west) rectangle (current page.north east);
  \path[fill=milcGradePrimary,rounded corners=16mm]
    ([xshift=.55cm,yshift=-.55cm]current page.north west) rectangle
    ([xshift=-.55cm,yshift=3.65cm]current page.south east);

  \node[anchor=north west,text=milcGradeText] at ([xshift=2.0cm,yshift=-1.85cm]current page.north west)
    {\sffamily\bfseries\fontsize{14}{16}\selectfont GRADO};
  \node[anchor=north west,text=milcGradeText,opacity=.82] at ([xshift=2.0cm,yshift=-2.75cm]current page.north west)
    {\sffamily\bfseries\fontsize{9}{11}\selectfont SERIE GUÍAS · TECNOLOGÍA E INFORMÁTICA};

  \begin{scope}[shift={([xshift=-3.05cm,yshift=-2.55cm]current page.north east)}]
    \fill[white,opacity=.80] (-.62,-.52) rectangle (.62,.52);
    \draw[milcGradeText,line width=1pt] (-.62,-.52) rectangle (.62,.52);
    \fill[milcGradeDark] (-.42,-.38) rectangle (-.22,.06);
    \fill[milcGradeAccent] (-.10,-.38) rectangle (.10,.24);
    \fill[milcGradePrimary!60!black] (.22,-.38) rectangle (.42,.38);
  \end{scope}

  \node[anchor=west,text=milcGradeText] at ([xshift=1.9cm,yshift=-12.85cm]current page.north west)
    {\sffamily\bfseries\fontsize{124}{124}\selectfont 6};
  \draw[milcGradeText,line width=1.7pt]
    ([xshift=4.52cm,yshift=-11.68cm]current page.north west) circle (.13cm);

  \node[anchor=west,text=milcGradeText,text width=8.7cm,align=left] at ([xshift=10.35cm,yshift=-11.6cm]current page.north west)
    {
     {\sffamily\bfseries\fontsize{19}{22}\selectfont Bienvenida ---\\cómo aprenderemos\\Tecnología\\este año}\\[4mm]
     {\sffamily\bfseries\fontsize{23}{26}\selectfont Guía 1-6-TIC}\\[2mm]
     {\sffamily\fontsize{13}{16}\selectfont Período 1 · Comunicación e identidad digital}\\[5mm]
     {\sffamily\bfseries\fontsize{11}{13}\selectfont Institución Educativa Sor María Juliana}\\[1mm]
     {\sffamily\fontsize{10}{12}\selectfont Autor: Dr. Álvaro Cárdenas Orozco}
    };

  \path[fill=milcGradeSoft,rounded corners=7mm]
    ([xshift=1.35cm,yshift=.85cm]current page.south west) rectangle
    ([xshift=-1.35cm,yshift=4.25cm]current page.south east);
  \draw[milcGradeDark,line width=.65pt,rounded corners=7mm]
    ([xshift=1.35cm,yshift=.85cm]current page.south west) rectangle
    ([xshift=-1.35cm,yshift=4.25cm]current page.south east);
  \node[anchor=center,text=milcNegro,text width=17.0cm,align=center] at ([yshift=2.55cm]current page.south)
    {\sffamily\fontsize{10}{12}\selectfont Metodología MILC · Apertura ancestral · Pensamiento computacional · Triángulo Dussel-Estoicismo-Floridi\\
    \textcolor{milcGradeDark}{\bfseries Producto:} Tu \textbf{cuaderno inaugural de Tecnología} listo para todo el año: portada con tu nombre + grado + año + materia (formato libre pero claro); índice de las \textbf{10 sesiones del periodo 1} (vas a llenar los títulos según la ruta que conozcas hoy); \textbf{compromiso firmado} de 5 acuerdos para trabajar bien este año (asistencia, cuaderno, respeto digital, esfuerzo, ayuda mutua). Más \textbf{en el cuaderno}: una pregunta que te llevas para investigar durante el periodo.};
\end{tikzpicture}
\mbox{}
\newpage

\section*{Apertura: Bienvenida a Tecnología --- cómo aprenderemos este año}

\begin{softbox}{milcGradeDark}{milcGradeSoft}{Saber ancestral}
\textbf{En los talleres de Cartago, recibir a un aprendiz nuevo era ritual.} Cuando un muchacho del campo llegaba donde el relojero don Lucho, donde la costurera doña Carmen o donde el panadero don Aurelio, \textbf{el primer día no se trabajaba}. Se conversaba. El maestro le mostraba el taller, le explicaba las herramientas, le decía qué se esperaba de él, qué horario, qué ritmo. \emph{``Aprenda primero a estar en el taller; después aprenderá el oficio''}, decía don Lucho a sus nuevos aprendices. Sin ese día de \textbf{bienvenida y orientación}, el aprendiz se sentía perdido y se iba la primera semana. Con ese día, se quedaba 5 años aprendiendo. \textbf{Nuestra sala de Tecnología es ese taller. Tú eres el aprendiz nuevo}. Hoy no vamos a empezar con Word, ni con internet, ni con hardware. Hoy nos conocemos, vemos el taller, hacemos el pacto. La próxima clase ya empieza el oficio.
\end{softbox}

\begin{softbox}{milcTurquesa}{milcGris}{Saber contemporáneo}
Este año en \textbf{Tecnología e Informática} vas a recorrer 3 grandes territorios. \textbf{Periodo 1 (Comunicación e identidad digital):} cómo presentarte en internet, cómo usar correo, cómo cuidarte de riesgos. \textbf{Periodo 2 (Hardware y software):} cómo funciona el computador por dentro y por fuera. \textbf{Periodo 3 (Procesador de texto e Internet):} cómo escribir documentos profesionales y cómo buscar información con criterio. En cada periodo tendrás \textbf{10 sesiones} (la 1 es esta apertura y la 10 cosecha lo aprendido) más una sesión 11 de cierre formal. Cada sesión sigue una estructura MILC: \textbf{saber ancestral} (cómo se hacía antes en Cartago), \textbf{saber contemporáneo} (lo técnico), \textbf{3 actividades} (escuchar, explicar, hacer), y \textbf{triángulo de pensamiento}. Tu cuaderno es la herramienta principal. El computador es el laboratorio. Aprender bien aquí te abre puertas para toda tu vida adulta.
\end{softbox}

\begin{softbox}{milcMostaza}{milcGris}{Pregunta-puente}
\textit{Si tuvieras que explicarle a tu mamá \emph{en qué consiste el grado 6° en Tecnología}, ¿\textbf{cómo lo harías} con lo que sabes hasta ahora? ¿Tienes una idea, o esperas que el año te lo muestre poco a poco?}
\end{softbox}

\begin{softbox}{milcVerde}{milcGris}{Saber y saber hacer}
Al terminar la clase: (1) podrás \textbf{identificar} la ruta del periodo 1 (los 10 temas); (2) sabrás \textbf{explicar} qué es el método MILC con tus palabras; (3) podrás \textbf{aplicar} las 5 reglas del cuaderno y del aula; (4) habrás \textbf{firmado} tu compromiso del año.
\end{softbox}

\small
\begin{tabularx}{\textwidth}{>{\bfseries}p{3.0cm}Y}
\rowcolor{milcGradePrimary!12}Autor & Dr. Álvaro Cárdenas Orozco\\
DBA & Reconoce principios y conceptos propios de la tecnología, relacionando artefactos con su utilización segura (MEN, T\&I 6°)\\
Referentes & Saber ancestral del pregonero · Pensamiento computacional inicial · Cuidado de la palabra\\
Producto final & Tu \textbf{cuaderno inaugural de Tecnología} listo para todo el año: portada con tu nombre + grado + año + materia (formato libre pero claro); índice de las \textbf{10 sesiones del periodo 1} (vas a llenar los títulos según la ruta que conozcas hoy); \textbf{compromiso firmado} de 5 acuerdos para trabajar bien este año (asistencia, cuaderno, respeto digital, esfuerzo, ayuda mutua). Más \textbf{en el cuaderno}: una pregunta que te llevas para investigar durante el periodo.\\
\end{tabularx}
\normalsize

\puente{Hoy haces 4 cosas: te presentas al curso, descubres la ruta del periodo, firmas el compromiso, y dejas el cuaderno listo.}

\begin{titlebox}{milcGradeDark}
Ruta MILC: así vamos a aprender
\end{titlebox}
\begin{tabularx}{\textwidth}{>{\centering\arraybackslash}X>{\centering\arraybackslash}X>{\centering\arraybackslash}X>{\centering\arraybackslash}X}
\phasecard{1}{milcVerde}{milcGris}{Escuta\\[-1mm]\small Observo, escucho y pregunto.} &
\phasecard{2}{milcTurquesa}{milcGris}{Sistematización\\[-1mm]\small Pensamiento computacional.} &
\phasecard{3}{milcMagenta}{milcGris}{Praxis\\[-1mm]\small Construyo y aplico.} &
\phasecard{4}{milcVino}{milcGris}{Evaluación\\[-1mm]\small Reflexión liberadora.}
\end{tabularx}


\puente{Empezamos con un círculo de la palabra: cada uno dice qué espera del año.}

\begin{titlebox}{milcVerde}
Fase 1 · Escuta
\end{titlebox}

\begin{softbox}{milcVerde}{milcGris}{Escena para escuchar}
\textbf{Actividad 1 · IDENTIFICA --- Círculo de la palabra} (10 min · individual + grupo). Sentados en U o en círculo, cada estudiante responde 3 preguntas en voz alta (1 frase corta cada una): \emph{(1) ¿Cómo te llamas y de qué barrio vienes?} \emph{(2) ¿Qué esperas de la materia de Tecnología este año?} \emph{(3) ¿Qué te gustaría aprender a hacer en el computador?}. El profe escribe en el tablero \textbf{3 palabras clave} que se repitan entre las respuestas (ejemplo: \emph{``jugar''}, \emph{``Word''}, \emph{``internet''}). Esas palabras se vuelven tu radar del año. \textbf{En el cuaderno:} apunta tus 3 respuestas. Después escribe las 3 palabras clave del grupo.
\end{softbox}

\checkbox Respondí las 3 preguntas en voz alta.\\
\checkbox Escuché las respuestas de los demás.\\
\checkbox Anoté en el cuaderno mis respuestas y las 3 palabras del grupo.

\begin{infoband}{milcVerde}
\textbf{Lo que va al cuaderno (Actividad 1 · IDENTIFICA).} Título: «Actividad 1 --- Círculo de la palabra». \textbf{Formato:} mis 3 respuestas en bloques + las 3 palabras clave del grupo. \textbf{Extensión:} 4 líneas. \textbf{Sabes que terminaste cuando} tienes registro escrito de quién eres y qué esperas, no solo en la cabeza.
\end{infoband}

\puente{Después conoces la ruta del periodo 1 (los 10 temas) y las reglas del trabajo.}

\begin{titlebox}{milcTurquesa}
Fase 2 · Sistematización con pensamiento computacional
\end{titlebox}

\begin{softbox}{milcTurquesa}{milcGris}{Cuatro pilares aplicados al tema}
Ahora conoce \textbf{la ruta del periodo 1} y las reglas del trabajo. Esto convierte la materia de \emph{misterio} en \emph{mapa}.
\end{softbox}

\small
\begin{tabularx}{\textwidth}{>{\bfseries\RaggedRight\arraybackslash}p{3.6cm}Y}
\rowcolor{milcTurquesa!90}\color{white}Pilar & \color{white}\bfseries Aplicación al tema\\
1. Descomposición & \textbf{La ruta del periodo 1 --- los 10 temas.} \emph{(S1) Hoy: Bienvenida.} \emph{(S2) Mi identidad digital.} \emph{(S3) Mi correo institucional.} \emph{(S4) Netiqueta: 10 acuerdos.} \emph{(S5) Historia de los medios.} \emph{(S6) Mi huella digital.} \emph{(S7) Privacidad y datos personales.} \emph{(S8) Riesgos digitales.} \emph{(S9) Comunicación responsable.} \emph{(S10) Mi presentación digital.} \emph{(S11) Cosecha y evaluación.} Cada una de S2 a S10 es \textbf{una sesión de 60-90 minutos} con un producto en el cuaderno. Si cumples las 10, sales del periodo 1 con la base sólida.\\
2. Reconocimiento de patrones & \textbf{El método MILC explicado simple.} Cada sesión empieza con un \textbf{saber ancestral} (cómo se hacía algo en Cartago antes), pasa por \textbf{3 actividades} (una de identificar, una de explicar, una de hacer), cierra con un \textbf{triángulo de pensamiento} (3 voces filosóficas: Dussel, Marco Aurelio o un estoico, y Floridi). \emph{No es decoración:} cada parte tiene un propósito. La ancestral te conecta con tu raíz; las actividades te entrenan en el oficio; el triángulo te entrena a pensar con voces que no son las tuyas. Al final, eres más libre y más sabio.\\
3. Abstracción & \textbf{Las 5 reglas del cuaderno y del aula.} \textbf{(1) Cuaderno propio:} 1 cuaderno solo para Tecnología, con cada sesión en página nueva, encabezado fijo. \textbf{(2) Asistencia:} las sesiones se construyen unas sobre otras. Si faltas una, te toca recuperarla solo. \textbf{(3) Respeto digital:} dentro y fuera del aula. Las reglas que aprenderás (netiqueta, privacidad) las aplicas también en tu día. \textbf{(4) Esfuerzo proporcional al objetivo:} no se pide perfección, se pide intento honesto. \textbf{(5) Ayuda mutua:} si terminas, ayuda al de al lado. Si no entiendes, pregunta sin pena. Aquí no se compite; se aprende juntos.\\
4. Algoritmo conceptual & \textbf{Encabezado fijo del cuaderno (cópialo en cada sesión nueva).} En la primera línea de cada sesión escribe: \emph{``Guía: 6° · Periodo [N] · Guía [N]''}. Segunda línea: \emph{``Tema: [título de la guía]''}. Tercera línea: \emph{``Fecha: \_\_/\_\_/2026''}. Esto te entrena en disciplina y le permite a tu yo de fin de año encontrar cualquier sesión en 30 segundos.\\
\end{tabularx}
\normalsize

\begin{drawbox}{milcTurquesa}{Ruta P1 + método MILC + 5 reglas}
\textbf{Ruta P1:} S1 (Bienvenida) → S2 (Identidad digital) → S3 (Correo) → S4 (Netiqueta) → S5 (Historia medios) → S6 (Huella) → S7 (Privacidad) → S8 (Riesgos) → S9 (Comunicación) → S10 (Presentación). \\[1mm]
\textbf{Método MILC:} saber ancestral + 3 actividades + triángulo. \\[1mm]
\textbf{5 reglas:} cuaderno propio · asistencia · respeto digital · esfuerzo honesto · ayuda mutua.
\end{drawbox}

\begin{softbox}{milcMostaza}{milcGris}{Errores comunes y su corrección}
\textbf{Errores frecuentes del primer día que te cuestan todo el año.} (1) \emph{No traer cuaderno propio o usar uno compartido con otra materia} --- todo se desorganiza rápido. (2) \emph{Faltar a la S1 pensando ``es solo bienvenida''} --- te pierdes la ruta y el pacto, y te cuesta meses ponerte al día. (3) \emph{Empezar sin compromiso firmado} --- no hay claridad de lo que se espera de ti. (4) \emph{Pensar que la materia es solo ``computador''} --- buena parte del oficio se hace en cuaderno físico. El computador es laboratorio, no televisor. (5) \emph{No preguntar cuando algo no está claro} --- las dudas iniciales se acumulan en confusiones grandes.
\end{softbox}

\begin{infoband}{milcTurquesa}
\textbf{Lo que va al cuaderno (Actividad 2 · EXPLICA).} Título: «Actividad 2 --- La ruta del periodo y mi compromiso». \textbf{Formato:} 2 bloques. Bloque A: lista numerada de los 10 temas del periodo. Bloque B: las 5 reglas con marca SÍ/NO si me comprometo con cada una. \textbf{Extensión:} 15 líneas. \textbf{Sabes que terminaste cuando} podrías decirle a tu mamá los 10 temas que vas a ver este periodo.
\end{infoband}

\puente{Con todo claro, organizas tu cuaderno y firmas el compromiso.}

\begin{titlebox}{milcMagenta}
Fase 3 · Praxis con pensamiento computacional
\end{titlebox}

\begin{softbox}{milcMagenta}{milcGris}{Construyendo el producto con los cuatro pilares}
\textbf{Actividad 3 · APLICA --- Tu cuaderno inaugural + compromiso firmado} (30 min · individual). Vas a dejar tu cuaderno LISTO para empezar el año. Es trabajo manual: portada, índice y compromiso, todo a mano.
\end{softbox}

\small
\begin{tabularx}{\textwidth}{>{\bfseries\RaggedRight\arraybackslash}p{3.6cm}Y}
\rowcolor{milcMagenta!90}\color{white}Pilar & \color{white}\bfseries Aplicación al producto\\
Descomposición & \textbf{Paso 1 --- Portada del cuaderno.} En la \textbf{primera hoja} (la de adentro de la tapa o la primera página) escribe en grande: \textbf{``Cuaderno de Tecnología e Informática''}. Debajo: tu nombre completo, grado y letra (6A, 6B, 6C, etc.), I.E. Sor María Juliana, año 2026. Decora la portada como tú quieras (dibujos, colores), pero el nombre y el grado deben ser legibles a 1 metro.\\
Patrones & \textbf{Paso 2 --- Índice del periodo 1.} En la \textbf{segunda hoja} escribe \textbf{``Índice --- Periodo 1''}. Debajo, lista numerada de las 10 sesiones del periodo (S1 a S10) con sus títulos cortos. Deja \textbf{una columna a la derecha vacía} con encabezado \emph{``Fecha de la sesión''} --- la llenarás cada vez que veas una sesión. Esto te da una vista del periodo completo.\\
Abstracción & \textbf{Paso 3 --- Compromiso del año.} En la \textbf{tercera hoja} escribe el título \textbf{``Mi compromiso del año en Tecnología''}. Debajo, escribe las 5 reglas (cuaderno propio, asistencia, respeto digital, esfuerzo honesto, ayuda mutua) y al lado de cada una, una palomita ✓ si te comprometes a cumplirla. Si dudas con alguna, conversa con el profe antes de firmar.\\
Algoritmo de armado & \textbf{Paso 4 --- Pregunta-radar.} Debajo del compromiso, escribe: \textbf{``Pregunta que me llevo para investigar este periodo''}. Después escribe \emph{1 pregunta tuya} que esperarías que se respondiera durante el periodo 1. Ejemplo: \emph{``¿Cómo me protejo de adultos extraños en internet?''}, o \emph{``¿De dónde sale realmente WhatsApp?''}. Esta pregunta te orienta durante las 10 sesiones.\\
\end{tabularx}
\normalsize

\begin{drawbox}{milcMagenta}{Antes de salir de la sala}
\checkbox Mi cuaderno tiene portada con nombre, grado, año, materia. \\[1mm]
\checkbox Hay un índice del periodo 1 con los 10 temas. \\[1mm]
\checkbox Las 5 reglas del compromiso están firmadas. \\[1mm]
\checkbox Tengo una pregunta-radar para el periodo. \\[1mm]
\checkbox La sesión 1 tiene su encabezado y las 3 actividades registradas. \\[1mm]
\checkbox Firmé con nombre y fecha al final.
\end{drawbox}

\begin{softbox}{milcMagenta}{milcGris}{Plantilla de trabajo}
\textbf{Estructura del cuaderno inaugural.} \textit{Hoja 1:} portada (nombre + grado + año + materia). \textit{Hoja 2:} índice P1 (10 temas + columna de fecha). \textit{Hoja 3:} compromiso (5 reglas con palomita + pregunta-radar). \textit{Hoja 4+:} sesión 1 con encabezado fijo + 3 actividades + firma.
\end{softbox}

\begin{infoband}{milcMagenta}
\textbf{Lo que va al cuaderno (Actividad 3 · APLICA).} Título: «Actividad 3 --- Cuaderno inaugural firmado». \textbf{Formato:} portada + índice + compromiso + pregunta + sesión 1 registrada. \textbf{Extensión:} 3-4 hojas. \textbf{Sabes que terminaste cuando} tu cuaderno está listo para empezar la próxima clase sin perder tiempo organizando.
\end{infoband}

\puente{El producto es tu cuaderno inaugural + el compromiso firmado.}

\begin{titlebox}{milcMostaza}
Producto de la sesión
\end{titlebox}
\textbf{Cuaderno inaugural de Tecnología · portada + índice + compromiso firmado}.
\begin{softbox}{milcMostaza}{milcGris}{Criterios de calidad}
(1) La portada tiene nombre completo, grado, año y materia legibles. (2) El índice lista las 10 sesiones del periodo 1 con columna de fecha. (3) Las 5 reglas del compromiso tienen palomita o firma. (4) Hay 1 pregunta-radar personal escrita. (5) La S1 tiene su encabezado fijo y las 3 actividades registradas.
\end{softbox}

\puente{Antes de salir, miras tu cuaderno: ¿está listo para empezar la próxima clase sin perder tiempo?}

\begin{titlebox}{milcVino}
Fase 4 · Evaluación liberadora — cinco dimensiones
\end{titlebox}

\begin{softbox}{milcVino}{milcGris}{Sentido de la evaluación}
Evaluar no es solo revisar una nota. Es reconocer qué aprendí, qué cambió en mí, qué emociones manejé, cómo me relaciono con la ciudad y la comunidad, y qué le dejaré a quien viene detrás.
\end{softbox}

\textbf{Mi reflexión liberadora en cinco dimensiones.} Completa cada frase iniciada:

\small
\begin{tabularx}{\textwidth}{>{\bfseries\RaggedRight\arraybackslash}p{3.4cm}Y>{\RaggedRight\arraybackslash}p{4.6cm}}
\rowcolor{milcVino}\color{white}Dimensión & \color{white}\bfseries Mi respuesta (frame starter) & \color{white}\bfseries Algo que descubrí\\
1. Personal & Lo que más cambió en mí fue \linead{6.5cm} & \linead{4.2cm}\\
2. Emocional & La emoción más fuerte fue \linead{2.5cm} y la regulé \linead{3cm} & \linead{4.2cm}\\
3. Ciudadana & En democracia esto importa porque \linead{6cm} & \linead{4.2cm}\\
4. Local (Cartago) & En mi barrio sería útil para \linead{6.5cm} & \linead{4.2cm}\\
5. Intergeneracional & A mi abuela/abuelo le diría \linead{6.5cm} & \linead{4.2cm}\\
\end{tabularx}
\normalsize

\begin{infoband}{milcVino}
\textbf{Para la próxima sustentación.} Vuelve a esta tabla en seis meses. Verás que las respuestas cambian — no porque lo aprendido cambie, sino porque tú cambias. La evaluación liberadora es una conversación contigo a través del tiempo.
\end{infoband}


\puente{Cerramos con 3 voces que te ayudan a entender por qué llegar bien es la mitad del camino.}

\begin{titlebox}{milcGradeDark}
Triángulo de pensamiento · cierre filosófico
\end{titlebox}

\begin{softbox}{milcVino}{milcGris}{Dussel · lente del nosotros}
\textit{````Llegar a un aula con cuaderno propio, nombre escrito y compromiso firmado, es declararse persona que aprende.''''}

Cuando ponees tu nombre en la portada y firmas el compromiso, no estás haciendo \emph{``trámite''}. Estás declarando: \emph{``yo soy estudiante de Tecnología este año, este es mi cuaderno, esta es mi voz''}. \textbf{Esa declaración es la primera forma de dignidad académica}. Los muchachos del taller de don Lucho hacían lo mismo cuando firmaban el contrato de aprendices. Tú hoy firmas el tuyo.

\textbf{¿Llegué hoy como persona que aprende, o como persona que pasa por la sala?}
\end{softbox}
\linead{\linewidth}\\[1mm]
\linead{\linewidth}

\begin{softbox}{milcMostaza}{milcGris}{Estoicismo · Marco Aurelio (emperador que escribió sobre cómo empezar cada año) · cuidado interior}
\textit{````Cada principio importa. El que empieza bien camina largo. El que empieza distraído, llega tarde toda la temporada.''''}

Marco Aurelio sabía que \textbf{el primer día marca el tono de los siguientes}. Si hoy llegas con cuaderno listo, atención puesta, compromiso firmado, has aprovechado bien la S1. Si llegas distraído, te demorará semanas ponerte al día. \emph{El primer día NO se recupera}. Por eso es la sesión más importante del periodo, aunque parezca solo bienvenida.

\textbf{¿Le di a hoy la importancia que tiene, o lo traté como un día cualquiera?}
\end{softbox}
\linead{\linewidth}\\[1mm]
\linead{\linewidth}

\begin{softbox}{milcTurquesa}{milcGris}{Floridi · lente de la infoesfera}
\textit{````Tecnología no se aprende sola. Se aprende con maestro, con compañeros, con cuaderno, con ritmo. La primera condición es haberlo entendido.''''}

Algunos estudiantes creen que \emph{``aprender computadores se hace solo viendo YouTube''}. Falso. \textbf{Se aprende con todo el ecosistema}: maestro que orienta, compañeros que dan retroalimentación, cuaderno que estructura, ritmo semanal que asienta. La S1 de hoy te enseñó eso: \emph{el oficio se aprende juntos, no en solitario}. Por eso la asistencia, el cuaderno físico y el compromiso son esenciales.

\textbf{¿Voy a tratar este año como solitario o como aprendizaje en comunidad?}
\end{softbox}
\linead{\linewidth}\\[1mm]
\linead{\linewidth}

\begin{titlebox}{milcVino}
Compromiso final
\end{titlebox}

\small
\begin{tabularx}{\textwidth}{>{\RaggedRight\arraybackslash}p{3.0cm}YYY}
\rowcolor{milcVino}\color{white}\bfseries Criterio & \color{white}\bfseries Lo logré & \color{white}\bfseries En proceso & \color{white}\bfseries Necesito apoyo\\
Comprensión del tema & Explico ideas centrales con mis palabras. & Reconozco ideas pero falta precisión. & Necesito repasar conceptos base.\\
Pensamiento computacional & Apliqué los 4 pilares al diseñar mi producto. & Apliqué algunos pilares. & Necesito releer la fase 2.\\
Aplicación y producto & Mi producto cumple los criterios y se puede revisar. & Producto funcional con ajustes pendientes. & Aún en construcción.\\
Reflexión 5 dimensiones & Respondí cada dimensión con honestidad. & Respondí algunas con detalle. & Mis respuestas son muy generales.\\
Triángulo de pensamiento & Conecté las 3 voces con mi trabajo real. & Conecté al menos una. & No relaciono las voces con mi proceso.\\
\end{tabularx}
\normalsize

\textbf{Hoy entendí que:}\\[1mm]
\linead{\linewidth}\\[1mm]
\linead{\linewidth}

\textbf{Mi compromiso para la próxima sesión es:}\\[1mm]
\linead{\linewidth}\\[1mm]
\linead{\linewidth}

\begin{softbox}{milcMostaza}{milcGris}{Cita de despedida · Marco Aurelio}
\textit{``El obstáculo en el camino se vuelve el camino. Lo que estorba al camino se vuelve camino.''}\\[1mm]
Cada error de hoy, bien observado, se convierte en pieza del camino que pisarás mañana.
\end{softbox}

\small
\begin{tabularx}{\textwidth}{>{\bfseries}p{3.6cm}Y}
\rowcolor{milcGradeDark!12}Nombre completo: & \linead{10cm}\\
Fecha: & \linead{4cm} \quad Curso/grupo: \linead{4cm}\\
Mi firma: & \linead{9cm}\\
\end{tabularx}
\normalsize

\begin{infoband}{milcGradeDark}
\textbf{Para la próxima sesión.} Trae tu \textbf{cuaderno listo} (portada + índice + compromiso firmado) y tu \textbf{lápiz/lapicero}. Mantén la pregunta-radar en la cabeza. La próxima clase empezamos el oficio: \emph{Mi identidad digital --- cómo me presento en la red sin perderme}. Vas a crear tu primera tarjeta de identidad digital. Hoy te recibiste como aprendiz; la próxima clase entras al taller.
\end{infoband}

\end{document}
